\documentclass[14pt,dvipdfmx]{jsarticle}
\input{lecture-preamble.tex}

% 例題3: y=-x^2+4x と点 A(1,3) における接線 y=2x+1
\newcommand{\exthreefig}{%
  \begin{tikzpicture}[x=0.72cm, y=0.46cm, line cap=round, line join=round]
    \draw[-{Stealth}] (-0.55,0) -- (4.85,0) node[right, font=\footnotesize] {$x$};
    \draw[-{Stealth}] (0,-0.85) -- (0,5.15) node[above, font=\footnotesize] {$y$};
    \node[below left, inner sep=1pt, font=\footnotesize] at (0,0) {O};
    \draw[thick, domain=0:4, samples=50] plot (\x, {-\x*\x+4*\x});
    \draw[red!75!black, thick] (-0.25,0.5) -- (1.9,4.8);
    \node[red!75!black, font=\footnotesize] at (0.12,1.95) {$\ell$};
    \draw[dashed, gray] (1,3) -- (1,0);
    \fill (1,3) circle (1.7pt);
    \node[above left, font=\footnotesize, inner sep=1.5pt] at (1,3) {A};
    \node[below, font=\footnotesize] at (1,0) {$1$};
    \node[below, font=\footnotesize] at (4,0) {$4$};
    \node[left, font=\footnotesize] at (0,3) {$3$};
  \end{tikzpicture}%
}

% 応用例題1: y=x^2+3 に点 C(1,0) から引いた2本の接線
\newcommand{\extangentfig}{%
  \begin{tikzpicture}[x=0.52cm, y=0.175cm, line cap=round, line join=round]
    \draw[-{Stealth}] (-2.15,0) -- (4.15,0) node[right, font=\footnotesize] {$x$};
    \draw[-{Stealth}] (0,-2.2) -- (0,14.6) node[above, font=\footnotesize] {$y$};
    \node[below left, inner sep=1pt, font=\footnotesize] at (0,0) {O};
    \draw[thick, domain=-1.65:3.45, samples=50] plot (\x, {\x*\x+3});
    \draw[red!75!black, thick] (-1.55,5.1) -- (1,0);
    \draw[red!75!black, thick] (1,0) -- (3.2,13.2);
    \fill (-1,4) circle (1.7pt);
    \fill (3,12) circle (1.7pt);
    \fill (1,0) circle (1.9pt);
    \node[anchor=north, font=\footnotesize, inner sep=1pt] at (1.15,-0.55) {C};
  \end{tikzpicture}%
}

\begin{document}

\lectureheader[6]{数学Ⅱ}{6}{微分法と積分法}{1}{微分係数と導関数}{3-AB}%
  {接線の方程式, グラフ上にない点から引いた接線}
  {188,189}{接線の方程式を求め, グラフの外の点からの接線も考えよう}

% ============ 1ページ目（表・左） ============
\vspace{-2.0zw}%
\heading{〜前時の復習〜}
\vspace{-0.55em}
\begin{reviewbox}{接線の傾き}
  \begin{enumerate}[label=\textbf{\arabic*}, leftmargin=1.5zw, itemsep=0.15em]
    \item グラフ上の点 $(a,\ f(a))$ における接線の傾きは,
          微分係数 \fitblankbf{f'(a)} に等しい.
    \item $f'(a)$ は, 導関数 $f'(x)$ に
          $x=\fitblankbf{a}$ を代入して求める.
  \end{enumerate}
\end{reviewbox}

\vspace{0.15em}
\topic{A}{接線の方程式}
\vspace{-0.55em}
関数 $y=f(x)$ のグラフ上の点 $(a,\ f(a))$ における接線の傾きは $f'(a)$ である.
この直線の方程式を求めよう.

\vspace{-0.15em}
\begin{exampleprob}{3}
  関数 $y=-x^{2}+4x$ のグラフ上に点 $\mathrm{A}(1,\ 3)$ をとる.
  \begin{tasks}[label=(\arabic*), label-width=1.6zw, item-indent=2zw, after-item-skip=0.04em](1)
    \task 点 $\mathrm{A}$ における接線 $\ell$ の傾き $m$ を求めよ.
    \task 点 $\mathrm{A}$ における接線 $\ell$ の方程式を求めよ.
  \end{tasks}
\end{exampleprob}
\vspace{-0.55em}
\begin{answerbox}{(1)}[9.6]
  $f(x)=-x^{2}+4x$ とすると, $m=f'(1)$ である.\\[0.2em]
  $f'(x)=-2x+4$ であるから\\[0.15em]
  $m=f'(1)=-2\cdot 1+4=2$
  \boxanswer{$2$}
\end{answerbox}

\vfill
\nextpage

% ============ 2ページ目（表・右） ============
\vspace{-1.4zw}%
\noindent\hbox to \linewidth{%
  \begin{minipage}[c]{0.58\linewidth}
    \begin{answerbox}{(2)}[12.2]
      接線 $\ell$ は, 点 $\mathrm{A}(1,\ 3)$ を通り,\\
      傾きが $2$ の直線である.\\[0.2em]
      よって $y-3=2(x-1)$\\[0.15em]
      すなわち $y=2x+1$
      \boxanswer{$y=2x+1$}
    \end{answerbox}
  \end{minipage}%
  \hfill
  \begin{minipage}[c]{0.40\linewidth}
    \centering\exthreefig
  \end{minipage}%
}

\vspace{0.25em}
これは, 点 $(1,\ f(1))$ を通り傾き $f'(1)$ の直線なので
\par\vspace{0.15em}
\hspace*{2zw}$y-f(1)=\fitblankbf{f'(1)(x-1)}$
\par\vspace{0.35em}
一般に, 次のことがいえる.

\vspace{-0.35em}
\begin{formulabox}{グラフ上の点における接線の方程式}
  関数 $y=f(x)$ のグラフ上の点 $(a,\ f(a))$ における接線の方程式は
  \par\vspace{0.15em}
  \centering $y-f(a)=\fitblankbf{f'(a)(x-a)}$
\end{formulabox}

\vspace{-0.15em}
\begin{practiceprob}{13}
  関数 $y=2x^{2}-4x+3$ のグラフ上に点 $\mathrm{A}(2,\ 3)$ をとる.
  公式 6.1.1 を用いて, 次を求めよ.
  \begin{tasks}[label=(\arabic*), label-width=1.6zw, item-indent=2zw, after-item-skip=0.02em](1)
    \task 点 $\mathrm{A}$ における接線の傾き
    \task 点 $\mathrm{A}$ における接線の方程式
  \end{tasks}
\end{practiceprob}
\vspace{-0.7em}
\begin{answerbox}{(1)}[6.6]
  $f(x)=2x^{2}-4x+3$ とすると $f'(x)=4x-4$\\[0.1em]
  傾きは $f'(2)=4\cdot 2-4=4$
  \boxanswer{$4$}
\end{answerbox}
\vspace{-0.35em}
\begin{answerbox}{(2)}[7.4]
  $y-f(2)=f'(2)(x-2)$ より $y-3=4(x-2)$\\[0.15em]
  すなわち $y=4x-5$
  \boxanswer{$y=4x-5$}
\end{answerbox}

\vfill
\nextpage

% ============ 3ページ目（裏・左） ============
\vspace{-1.4zw}%
\topic{B}{グラフ上にない点から引いた接線}
\vspace{-0.7em}
公式 6.1.1 は, 接点 $(a,\ f(a))$ における接線の方程式である.
グラフ上にない点から接線を引くときは, 接点がわからない.
接点の $x$ 座標を \fitblankbf{a} とおき,
接線がその点を\fitblankbf{\text{通る}}ことから $a$ を求める.

\vspace{0.15em}
\noindent{\bfseries 応用例題\,1}\quad
関数 $y=x^{2}+3$ のグラフに点 $\mathrm{C}(1,\ 0)$ から引いた接線は $2$ 本ある.
この $2$ 本の接線の方程式を求めよ.
\par\vspace{0.35em}
\noindent\hbox to \linewidth{%
  \begin{minipage}[c]{0.56\linewidth}
    \textbf{考え方}\quad
    接点の $x$ 座標を $a$ として,
    点 $(a,\ a^{2}+3)$ における接線の方程式を求める.
    その接線が点 $\mathrm{C}$ を通る式から, $a$ が決まる.
  \end{minipage}%
  \hfill
  \begin{minipage}[c]{0.42\linewidth}
    \centering\extangentfig
  \end{minipage}%
}

\vspace{0.15em}
\begin{answerbox}{}[28]
  $y=x^{2}+3$ を微分すると $y'=2x$\\[0.2em]
  接点を $(a,\ a^{2}+3)$ とすると, 傾きは $2a$ であるから\\[0.12em]
  $y-(a^{2}+3)=2a(x-a)$\\[0.12em]
  すなわち $y=2ax-a^{2}+3$ ……①\\[0.28em]
  この直線が点 $\mathrm{C}(1,\ 0)$ を通るから\\[0.12em]
  $0=2a-a^{2}+3$ \quad よって $a^{2}-2a-3=0$\\[0.12em]
  $(a+1)(a-3)=0$ \quad したがって $a=-1,\ 3$\\[0.28em]
  ①より, $a=-1$ のとき $y=-2x+2$\\[0.1em]
  $a=3$ のとき $y=6x-6$
  \boxanswer{$y=-2x+2,\quad y=6x-6$}
\end{answerbox}

\vfill
\nextpage

% ============ 4ページ目（裏・右） ============
\vspace{-1.4zw}%
\textbf{補足}\quad
点 $\mathrm{C}(1,\ 0)$ を通る直線は, 傾きを $m$ とすると
$y=m(x-1)$ と表される.
応用例題1では $m=\fitblankbf{-2},\ \fitblankbf{6}$ である.

\vspace{0.05em}
\begin{practiceprob}{14}
  関数 $y=x^{2}-2x+4$ のグラフに原点 $\mathrm{O}$ から引いた接線は $2$ 本ある.
  この $2$ 本の接線の方程式を求めよ.
\end{practiceprob}
\vspace{-0.55em}
\begin{answerbox}{}[22.5]
  $y'=2x-2$\\[0.15em]
  接点を $(a,\ a^{2}-2a+4)$ とすると, 傾きは $2a-2$ であるから\\[0.1em]
  $y=(2a-2)(x-a)+a^{2}-2a+4$\\[0.12em]
  $=(2a-2)x-a^{2}+4$ ……①\\[0.22em]
  原点 $(0,\ 0)$ を通るから $0=-a^{2}+4$\\[0.1em]
  $a^{2}=4$ \quad よって $a=2,\ -2$\\[0.22em]
  ①より, $a=2$ のとき $y=2x$\\[0.08em]
  $a=-2$ のとき $y=-6x$
  \boxanswer{$y=2x,\quad y=-6x$}
\end{answerbox}

\vspace{0.15em}
\textbf{Point}\quad
接線が通る点が与えられたとき, その点が\fitblankbf{\text{接点}}かどうかで解法が異なる.
\par\vspace{0.12em}
\hspace*{2zw}例題3は, 与えられた点が接点である.\quad
応用例題1は, 与えられた点が接点でない.

\vfill
\begin{summarybox}
  \begin{enumerate}[label=\textbf{\arabic*}, leftmargin=1.5zw, itemsep=0.28em, topsep=0.12em]
    \item 点 $(a,\ f(a))$ における接線の方程式は
          \par\vspace{0.12em}
          {\renewcommand{\fitblankpad}{1.2zw}%
          \hspace*{1zw}\fitblankbf{y-f(a)=f'(a)(x-a)}}
          \par\vspace{0.2em}
    \item グラフ上にない点から引くときは, 接点の $x$ 座標を $a$ とおき,
          接線がその点を通ることから $a$ を求める.
  \end{enumerate}
\end{summarybox}

\vspace{0.05em}
\nexttime{関数の増減を, 導関数の符号から考えてみよう!}

\end{document}
